\subsection{Random Variables and Probability Mass Functions}

\begin{definition}{Random Variable}{random-variable}
  A random variable on a sample space $S$ is a function $X:S\to\R$.
\end{definition}

\begin{example}
  $S=\{THT\cdots,\ldots\}$, the sample space of tossing a coin $n$ times.
  For $\omega\in S$,
  \[
    X(\omega)=\#\text{ H's in }\omega
  \]
  is a random variable.

  Another random variable:
  \[
    X(\omega)=
    \begin{cases}
      \text{position of first H in }\omega, & \text{if there is an H}, \\
      n+1,                                  & \text{if no H}.
    \end{cases}
  \]
\end{example}

\begin{example}
  $S=\{(1,1),(1,2),\ldots,(6,6)\}$.
  For $\omega=(\omega_1,\omega_2)\in S$,
  \[
    X(\omega)=\omega_1+\omega_2.
  \]
\end{example}

\begin{example}[Largest Number Selected]
  20 numbered balls in an urn; choose 4 at random.
  \[
    S=\bigl\{\{1,2,3,4\},\{1,2,3,5\},\ldots,\{17,18,19,20\}\bigr\}.
  \]
  $X(\omega)=$ largest number selected.
  \begin{gather*}
    P(X=4)=\frac1{\binom{20}{4}},\qquad
    P(X=5)=\frac{\binom43}{\binom{20}{4}},\\
    P(X=6)=\frac{\binom53}{\binom{20}{4}},\qquad\ldots,\qquad
    P(X=k)=\frac{\binom{k-1}{3}}{\binom{20}{4}},\quad k=4,\ldots,20.
  \end{gather*}
  \[
    \sum_{k=4}^{20}P(X=k)=1\qquad\text{(all of sample space)}.
  \]
\end{example}

\begin{definition}{Probability Mass Function}{probability-mass-function}
  For a discrete random variable $X$,
  \[
    p(x)=P(X=x)
  \]
  is the probability mass function (pmf) of $X$.
\end{definition}

For a random variable, we have $E=\{\omega:X(\omega)=x\}\subseteq S$,
\[
  P(E)=P(\{\omega:X(\omega)=x\})=P(X=x)=p(x).
\]
Properties:
\[
  p(x)\ge0,\qquad\sum_{x:p(x)>0}p(x)=1.
\]
\[
  P(X\in E)=\sum_{x\in E}p(x).
\]

\begin{definition}{Cumulative Distribution Function}{cumulative-distribution-function}
  $F(x)=P(X\le x)$ is called the cumulative distribution function of $X$.
\end{definition}

\subsection{Discrete Distributions}

\begin{definition}{Uniform Distribution}{discrete-uniform-distribution}
  Uniform on $\{1,2,\ldots,n\}$:
  \[
    p(x)=P(X=x)=\frac1n,\qquad x=1,\ldots,n.
  \]
\end{definition}

\begin{example}[Toss Until Heads]
  Toss a coin until H appears.
  \[
    S=\{H,TH,TTH,\ldots\},\qquad P(H)=p\in(0,1).
  \]
  $X(\omega)=$ length of $\omega$.
  \[
    \{\omega:X(\omega)=k\}=\{X=k\}.
  \]
  \begin{align*}
    P(X=1) & =p,            \\
    P(X=2) & =(1-p)p,       \\
    P(X=3) & =(1-p)^2p,     \\
    P(X=k) & =(1-p)^{k-1}p.
  \end{align*}
\end{example}

\begin{definition}{Geometric Distribution}{geometric-distribution}
  Geometric pmf with parameter $p\in(0,1)$:
  \[
    P(X=x)=p(x)=(1-p)^{x-1}p,\qquad x=1,2,\ldots.
  \]
  First success on $x$.\\
  A random variable with geometric pmf is a geometric random variable.
\end{definition}

\begin{gather*}
  \sum_{k=1}^{\infty}p(1-p)^{k-1}\stackrel{?}{=}1,\\
  \text{WTS }p\sum_{k=1}^{\infty}(1-p)^{k-1}=1,\\
  \text{WTS }p\sum_{\ell=0}^{\infty}(1-p)^\ell=1.
\end{gather*}
Geometric series:
\[
  |r|<1\implies\sum_{k=0}^{\infty}r^k=\frac1{1-r}.
\]

\begin{example}[Number of Heads]
  Toss a coin $n$ times (tosses independent). $S$ is the sample space:
  \[
    S=\{HH\cdots H,HTH\cdots H,HHTH\cdots H,\ldots,TT\cdots T\}.
  \]
  An $\omega\in S$ is a sequence of length $n$,
  \[
    \omega=(\omega_1,\omega_2,\ldots,\omega_n),\qquad\omega_i\in\{H,T\}.
  \]
  \[
    X(\omega)=\#\{i:\omega_i=H\}.
  \]
  $P(X=k)$: probability that the number of H's is $k$.

  For $n=5$, $k=3$:
  \[
    \begin{array}{ccccc}
      H & H & H & T & T \\
      H & H & T & H & T \\
      \multicolumn{5}{c}{\cdots}
    \end{array}
  \]
  \[
    P(X=k)=\binom nk p^k(1-p)^{n-k}.
  \]
\end{example}

\begin{definition}{Binomial Distribution}{binomial-distribution}
  Binomial pmf:
  \[
    p(k)=\binom nk p^k(1-p)^{n-k},\qquad k=0,\ldots,n.
  \]
  Written as $\operatorname{Bi}(n,p)$, with success probability $p$, $n$ trials.
\end{definition}

\begin{definition}{Poisson Distribution}{poisson-distribution}
  $X$ is called a Poisson random variable with parameter $\lambda>0$ if
  \[
    P(X=x)=\frac{\lambda^x}{x!}e^{-\lambda}=p(x),\qquad x=0,1,2,\ldots.
  \]
\end{definition}

Note: for the Poisson pmf,
\[
  \sum_{m=0}^{\infty}\frac{\lambda^m}{m!}e^{-\lambda}
  =e^{-\lambda}\sum_{m=0}^{\infty}\frac{\lambda^m}{m!}
  =e^{-\lambda}e^\lambda=1.
\]

\subsubsection{Binomial to Poisson}

Note that
\begin{align*}
  \text{Poisson r.v.: }  & p(k)=\frac{\lambda^k}{k!}e^{-\lambda}, \\
  \text{Binomial r.v.: } & p(k)=\binom nk p^k(1-p)^{n-k}.
\end{align*}
For $Y\sim\operatorname{Bi}(n,p)$, take $p=\lambda/n$ and
$Y_n\sim\operatorname{Bi}(n,\lambda/n)$.
\begin{align*}
  P(Y_n=0) & =\left(1-\frac\lambda n\right)^n\to e^{-\lambda},           \\
  P(Y_n=1) & =n\frac\lambda n\left(1-\frac\lambda n\right)^{n-1}
  \to\lambda e^{-\lambda},                                               \\
  P(Y_n=2) & =\frac{n(n-1)}2\frac{\lambda^2}{n^2}
  \left(1-\frac\lambda n\right)^{n-2}\to\frac{\lambda^2}{2}e^{-\lambda}, \\
  P(Y_n=3) & =\frac{n(n-1)(n-2)}{3\cdot2\cdot1}\frac{\lambda^3}{n^3}
  \left(1-\frac\lambda n\right)^{n-3}\to\frac{\lambda^3}{3!}e^{-\lambda}.
\end{align*}
Poisson distribution.

\subsection{Expected Value}

\begin{definition}{Expected Value}{expected-value}
  For a random variable $X$ with pmf $p$ and $\sum_x|x|p(x)<\infty$,
  \[
    EX=\sum_x xp(x)
  \]
  is the expected value of $X$.
\end{definition}

\begin{example}[Uniform]
  \[
    EX=\sum_x xp(x)=\sum_{k=1}^n k\cdot\frac1n
    =\frac1n\cdot\frac{n(n+1)}2=\frac{n+1}2.
  \]
\end{example}

\begin{example}[Binomial]
  For $\operatorname{Bi}(n,p)$:
  \begin{align*}
    EX & =\sum_{k=0}^n kp(k)=\sum_{k=0}^n k\binom nk p^k(1-p)^{n-k}       \\
       & =\sum_{k=0}^n k\frac{n!}{k!(n-k)!}p^k(1-p)^{n-k}                 \\
       & =\sum_{k=1}^n\frac{n!}{(k-1)!(n-k)!}p^k(1-p)^{n-k}               \\
       & =np\sum_{k=1}^n\frac{(n-1)!}{(k-1)!(n-1-(k-1))!}
    p^{k-1}(1-p)^{n-1-(k-1)}                                              \\
       & =np\sum_{j=0}^{n-1}\frac{(n-1)!}{j!((n-1)-j)!}p^j(1-p)^{(n-1)-j} \\
       & =np\sum_{j=0}^{n-1}\binom{n-1}{j}p^j(1-p)^{n-1-j}                \\
       & =np[p+(1-p)]^{n-1}=np.
  \end{align*}
\end{example}

\begin{example}[Poisson]
  For $\operatorname{Po}(\lambda)$:
  \begin{align*}
    EX & =\sum_{k=0}^{\infty}k\frac{\lambda^k}{k!}e^{-\lambda}
    =e^{-\lambda}\sum_{k=0}^{\infty}k\frac{\lambda^k}{k!}                      \\
       & =e^{-\lambda}\sum_{k=1}^{\infty}\frac{\lambda^k}{(k-1)!}              \\
       & =\lambda e^{-\lambda}\sum_{j=0}^{\infty}\frac{\lambda^j}{j!}=\lambda.
  \end{align*}
\end{example}

Till now, only countable sets and discrete random variables.

\subsection{Coupon Collector}

\begin{example}[Collecting a Complete Set]
  Collect $r$ types of coupons; receive a reward for a complete set.
  Assume all types equally likely and successive collections independent.

  Let $T=$ number of collections until a complete set is collected.
  Apparently $P(T<r)=0$.

  $P(T=n)$ for $n\in\{r,r+1,r+2,\ldots\}$?
  \[
    \underbrace{\{T=n\}\cup\{T>n\}}_{\text{disjoint}}=\{T>n-1\}
    \qquad\text{(immediately obvious)}.
  \]
  \begin{gather*}
    P(\{T=n\}\cup\{T>n\})=P(T=n)+P(T>n)=P(T>n-1),\\
    P(T=n)=P(T>n-1)-P(T>n).
  \end{gather*}
  $P(T>n)$: at least one coupon type missing.

  For $j\in\{1,2,\ldots,r\}$, define
  \[
    A_j=\{\text{haven't got type }j\text{ by the }n\text{th collection}\}.
  \]
  \[
    \{T>n\}=\bigcup_{j=1}^r A_j.
  \]
  \begin{align*}
    P\left(\bigcup_{j=1}^r A_j\right)
     & =\sum_{j=1}^r P(A_j)-\sum_{i<j}P(A_iA_j) \\
     & \quad+\sum_{i<j<k}P(A_iA_jA_k)-\cdots
    +(-1)^{r+1}P(A_1\cdots A_r).
  \end{align*}
  \[
    P(A_j)=\left(\frac{r-1}{r}\right)^n,\qquad
    P(A_iA_j)=\left(\frac{r-2}{r}\right)^n\quad(i\ne j).
  \]
  \begin{align*}
    P\left(\bigcup_{j=1}^r A_j\right)
     & =\binom r1\left(\frac{r-1}{r}\right)^n
    -\binom r2\left(\frac{r-2}{r}\right)^n                \\
     & \quad+\binom r3\left(\frac{r-3}{r}\right)^n-\cdots
    +(-1)^r\binom r{r-1}\left(\frac1r\right)^n            \\
     & =\sum_{\ell=1}^{r-1}(-1)^{\ell+1}\binom r\ell
    \left(\frac{r-\ell}{r}\right)^n.
  \end{align*}
\end{example}
